\documentclass[10pt,a4paper]{amsart}
\usepackage[margin=19mm]{geometry}
\usepackage{amsmath,amssymb,mathtools}
\usepackage[scr=rsfs,cal=euler]{mathalfa}
\usepackage{xfrac}
\usepackage[unicode,hidelinks,bookmarks=false]{hyperref}
\newcommand{\ie}{i.e.\ }
\newcommand{\cf}{cf.\ }
\newcommand{\resp}{respectively\ }
\newcommand{\Subsection}{\subsection}
\providecommand{\nobreakdash}{\nobreak\hbox{-}}
\setlength{\emergencystretch}{2em}
\multlinegap0pt
\usepackage{bm}
\usepackage{bbm}
\usepackage{dsfont}
\usepackage{xspace}
\usepackage{cancel}
\usepackage{mathtools}
\usepackage{amsthm}
\makeatletter
\@ifundefined{assumption}{\newtheorem{assumption}{Assumption}}{}
\@ifundefined{proposition}{\newtheorem{proposition}{Proposition}}{}
\makeatother
\newcommand{\qvec}{\mathbf q}
\newcommand{\xvec}{\mathbf x}
\title[A factorisation-based regularised interior point method usin]{A factorisation-based regularised interior point method using the augmented system}
\author{Filippo Zanetti, Jacek Gondzio}
\date{Source: arXiv:2508.04370v1 (CC BY 4.0)}
\begin{document}
\maketitle

\noindent\emph{Source and licence.} A factorisation-based regularised interior point method using the augmented system. Version arXiv:2508.04370v1 (CC BY 4.0), distributed under the Creative Commons Attribution 4.0 International licence, as recorded in the arXiv licence metadata for this exact version and shown on the abstract page https://arxiv.org/abs/2508.04370v1. Full rights notice: licenses/arxiv-2508.04370-CC-BY-4.0.txt.\par\smallskip

\noindent\textit{Editorial scope.} Counted unit: the Proposition whose source label is \texttt{None} in the article (source0.tex lines 488--490, source label \texttt{None}), delivered together with the article's own complete original proof of it (source0.tex lines 491--514, one \texttt{proof} environment of 24 lines). No mathematics is written, repaired or re-derived: statement and proof are the article's own text line for line. Required context retained uncounted: assum:structure (lines 464--475); assum:quasi (lines 478--484). Every definition, display and label of those lines is retained unchanged. Omitted from this package and not redistributed: 1--463, 476--477, 485--487, 515--1191. Nothing inside a retained excerpt is omitted or altered: every retained body is delivered exactly as the article's lines are written, so no omission receipt of any kind accompanies this package. Citations: the retained excerpts carry no citation command at all, so no bibliography entry of the article is transcribed and no reference is redistributed. Reference adaptation: none was needed and none was made; every reference inside the delivered unit resolves inside it, so no unresolved reference or citation is produced. Printed slips kept literal: no printed word, symbol, hypothesis or number of the counted statement or of its proof was altered. Layout and class: the article's own class is \texttt{article} with packages including fontenc, lmodern, babel, inputenc, setspace, hyperref, microtype, quoting, amsmath, amssymb, comment, amsthm, subcaption, mathrsfs; the wrapper is the lane's pinned \texttt{amsart} editorial wrapper at 10pt on A4, which restates the theorem-like environments the retained text needs and adds the article's own macro definitions that the retained text needs (\texttt{\textbackslash qvec}, \texttt{\textbackslash xvec}), with extra wrapper packages (bm, bbm, dsfont, xspace, cancel, mathtools). The delivered numbering is the unit's own. Assets: no retained span contains a figure environment, a graphics-inclusion command, a PDF-page inclusion, a table or a TikZ picture; no image file is copied into this package, the package declares no asset dependency and no image file is redistributed with it. Licence: version arXiv:2508.04370v1 of this article is distributed under the Creative Commons Attribution 4.0 International licence, as recorded in the arXiv licence metadata for this exact version and shown on the abstract page https://arxiv.org/abs/2508.04370v1. A full rights notice is delivered at licenses/arxiv-2508.04370-CC-BY-4.0.txt. Characters: every retained excerpt is pure ASCII, so the delivered engine, XeLaTeX with its default Latin Modern text font, reports no missing character.
\begin{assumption}
\label{assum:structure}
The matrix has the following structure
\[
\begin{bmatrix}
p & \qvec_1^T & \qvec_2^T \\
\qvec_1 & M_{11} & M_{21}^T \\
\qvec_2 & M_{21} & M_{22}
\end{bmatrix}
\]
where $p<0$, $\qvec_1\in\mathbb R^{n_1}$, $\qvec_2\in\mathbb R^{n_2}$, $M_{11}\in\mathbb R^{n_1\times n_1}$, $M_{21}\in\mathbb R^{n_2\times n_1}$, $M_{22}\in\mathbb R^{n_2\times n_2}$, $M_{11}\prec 0$, $M_{22}\succ0$.
\end{assumption}

\begin{assumption}
\label{assum:quasi}
\[\begin{bmatrix}
p & \qvec_1^T \\ \qvec_1 & M_{11}
\end{bmatrix}\prec0.
\]
\end{assumption}

\begin{proposition}
If Assumptions \ref{assum:structure} and \ref{assum:quasi} hold, then, after eliminating pivot $p$, the remaining Schur complement is quasi-definite.
\end{proposition}

\begin{proof}
Eliminating pivot $p$ leaves a matrix of size $n_1+n_2$ to be processed in the next steps. The $(1,1)-$block of the remaining matrix is $M_{11}-\frac{1}{p}\qvec_1\qvec_1^T$; the $(2,2)-$ block reads $M_{22}-\frac{1}{p}\qvec_2\qvec_2^T$. Since $p<0$ and $M_{22}\succ0$, the $(2,2)-$block is strictly positive definite. 

Consider now Assumption~\ref{assum:quasi}, i.e.,
\[
\begin{bmatrix}
\alpha & \xvec^T
\end{bmatrix}
\begin{bmatrix}
p & \qvec_1^T \\ \qvec_1 & M_{11}
\end{bmatrix}
\begin{bmatrix}
\alpha \\ \xvec
\end{bmatrix}
<0
\]
for any $\alpha\in\mathbb R$, $\xvec\in\mathbb R^{n_1}$. This is equivalent to 
\begin{align}
0>\,\,& \alpha^2p+2\alpha \qvec_1^T\xvec +\xvec^TM_{11}\xvec \notag\\
=\,\,& \xvec^TM_{11}\xvec-\frac{1}{p}(\qvec_1^T\xvec)^2+\frac{1}{p}\Big[(\qvec_1^T\xvec)^2+2\alpha p \qvec_1^T\xvec+\alpha^2 p^2\Big] \notag\\
=\,\,& \xvec^T\Big(M_{11}-\frac{1}{p}\qvec_1\qvec_1^T\Big)\xvec+\frac{1}{p}\Big(\qvec_1^T\xvec+\alpha p\Big)^2 \notag.
\end{align}
Choosing $\alpha=-\qvec_1^T\xvec/p$ yields that $M_{11}-\frac{1}{p}\qvec_1\qvec_1^T$ is negative definite. Therefore, after eliminating pivot $p$, the remaining Schur complement has a negative definite $(1,1)-$block and a positive definite $(2,2)-$block.
\end{proof}

\end{document}
